\documentclass[uplatex,dvipdfmx]{jsarticle}
\usepackage{amsmath}
\usepackage{graphicx}
\usepackage{url}
\usepackage{listings}
\usepackage{xcolor}

\lstset{
  basicstyle=\ttfamily\small,
  breaklines=true,
  frame=single,
  columns=fullflexible,
  backgroundcolor=\color{gray!8}
}

\title{入力映像に基づく顔モデルテクスチャの色味補正}
\author{}
\date{}

\begin{document}
\maketitle

\section{目的}

本研究では，カメラ映像または入力動画上の顔と，重畳表示する3次元顔モデルとの見た目の差を小さくすることを目的とする．
特に，モデルの肌色が入力映像と異なる場合，モデルが浮いて見える問題がある．
そこで，入力映像から肌色の基準色を取得し，モデルテクスチャ全体の色味を補正する処理を実装した．

\section{従来の問題}

初期状態では，モデルテクスチャは生成時の画像の色をそのまま使用していた．
そのため，入力映像の照明条件やカメラ特性が変化すると，以下のような問題が生じた．

\begin{itemize}
  \item 入力映像の顔とモデルの肌色が一致しない
  \item モデルだけ明るい，または薄く見える
  \item 1点のみを基準にすると，影やハイライトの影響を受けやすい
  \item 補正を強くしすぎると，口や鼻などのテクスチャ情報が失われる
\end{itemize}

\section{提案手法}

本実装では，MediaPipe FaceMesh の複数ランドマークを肌色基準点として使用する．
単一のランドマークではなく複数点を用いることで，局所的な影，ハイライト，検出誤差の影響を抑える．

使用する肌色基準ランドマークは以下である．

\begin{lstlisting}[language=Python]
skin_landmarks = [10, 151, 9, 8, 168, 197, 6, 195, 5, 4]
\end{lstlisting}

各ランドマーク周辺の小領域からRGB値を取得し，それらをまとめた中央値を肌色の代表値とする．
中央値を用いることで，外れ値となる明るい画素や暗い画素の影響を小さくできる．

\section{補正式}

モデルテクスチャ側の肌色基準を $\mathbf{C}_{model}$，入力映像側の肌色基準を $\mathbf{C}_{input}$ とする．
また，モデルテクスチャの各画素のRGB値を $\mathbf{T}$，補正後のRGB値を $\mathbf{T}'$ とする．
実装では，以下の式で補正を行う．

\begin{equation}
  \mathbf{T}' = \mathbf{T} + (\mathbf{C}_{input} - \mathbf{C}_{model})
\end{equation}

この式の意味は，モデルテクスチャの基準肌色からの相対的な差分を保持しつつ，基準となる肌色だけを入力映像側へ移動することである．
そのため，口や鼻，陰影などのテクスチャ情報を残したまま，肌色を入力映像に近づけることができる．

\section{色補正方式の比較}

現在の実装では，モデルテクスチャの補正方式を切り替えられるようにしている．
補正方式は大きく分けて，RGB各チャンネルを直接補正する方法と，色空間を変換して輝度成分のみを補正する方法がある．

\subsection{RGB差分補正}

RGB差分補正では，モデル側の基準肌色 $\mathbf{C}_{model}$ と入力映像側の基準肌色 $\mathbf{C}_{input}$ の差分を，テクスチャ全体のRGB値へ直接加算する．

\begin{equation}
  \mathbf{T}' = \mathbf{T} + (\mathbf{C}_{input} - \mathbf{C}_{model})
\end{equation}

これは，R，G，B の各チャンネルを独立に移動する処理である．
例えば入力映像の肌がモデルより赤い場合はR成分が増え，青みが少ない場合はB成分が減る．
したがって，明るさだけでなく，赤み，黄み，青みといった色味も入力映像へ近づけられる．

\begin{lstlisting}[language=Python]
rgb = (rgb - source_rgb) + target_rgb
\end{lstlisting}

この方式の利点は，モデルの肌色が入力映像と大きく異なる場合でも，色味そのものを強く合わせられることである．
一方で，カメラ映像の一時的な色かぶりやランドマーク周辺の影響を受けると，モデル全体の色相が不自然に変化する可能性がある．

\subsection{YCrCb輝度補正}

YCrCb輝度補正では，RGB画像をYCrCb色空間へ変換し，輝度を表すY成分のみを補正する．
Cr成分とCb成分は色差成分であり，赤みや青みなどの色味を表す．
そのため，Yのみを変更すると，モデルテクスチャ本来の色味をできるだけ保ったまま，明るさだけを入力映像へ近づけることができる．

モデル側の基準輝度を $Y_{model}$，入力映像側の基準輝度を $Y_{input}$，テクスチャ画素の輝度を $Y_T$ とすると，補正は以下のようになる．

\begin{equation}
  Y'_T = Y_T + (Y_{input} - Y_{model})
\end{equation}

補正後は，変更した $Y'_T$ と元のCr，Cbを用いてRGBへ戻す．

\begin{lstlisting}[language=Python]
ycrcb = cv2.cvtColor(rgb, cv2.COLOR_RGB2YCrCb)
ycrcb[:, :, 0] += target_y - source_y
rgb = cv2.cvtColor(ycrcb, cv2.COLOR_YCrCb2RGB)
\end{lstlisting}

この方式は，モデルの色味自体は大きく変えたくないが，照明や露出の差によってモデルだけ暗い，または明るい場合に有効である．
RGB差分補正より色相の揺れが少ないため，フレーム間のちらつきも抑えやすい．
ただし，モデルの肌色と入力映像の肌色が根本的に異なる場合，明るさは合っても色味の差は残る．

\subsection{Lab輝度補正}

Lab輝度補正では，RGB画像をLab色空間へ変換し，明度を表すL成分のみを補正する．
Lab色空間では，Lが明るさ，aとbが色味を表す．
YCrCbと同様に，Lのみを変更することで，色味成分を保ったまま明るさを調整できる．

モデル側の基準明度を $L_{model}$，入力映像側の基準明度を $L_{input}$，テクスチャ画素の明度を $L_T$ とすると，補正は以下のようになる．

\begin{equation}
  L'_T = L_T + (L_{input} - L_{model})
\end{equation}

補正後は，変更した $L'_T$ と元のa，bを用いてRGBへ戻す．

\begin{lstlisting}[language=Python]
lab = cv2.cvtColor(rgb, cv2.COLOR_RGB2LAB)
lab[:, :, 0] += target_l - source_l
rgb = cv2.cvtColor(lab, cv2.COLOR_LAB2RGB)
\end{lstlisting}

Labは人間の知覚に近い形で明るさと色味を分離するため，輝度補正の考え方としては自然である．
一方で，OpenCV上のLab変換では値域が8bit表現に丸められるため，変換と逆変換の過程でRGB値がわずかに変化する場合がある．

\subsection{方式ごとの使い分け}

各方式の特徴をまとめると以下のようになる．

\begin{itemize}
  \item RGB差分補正は，肌色そのものを入力映像へ強く合わせたい場合に有効である
  \item YCrCb輝度補正は，モデル本来の色味を保ちつつ，カメラ露出や照明差だけを補正したい場合に有効である
  \item Lab輝度補正は，知覚的な明るさを基準にして，色味を保ったまま自然な明度補正を行いたい場合に有効である
\end{itemize}

実装上は，以下の設定で補正方式を切り替える．

\begin{lstlisting}[language=Python]
model_color_match_mode = 'rgb'
model_color_match_mode = 'ycrcb_luminance'
model_color_match_mode = 'lab_luminance'
\end{lstlisting}

現在の確認では，輝度補正の動作を確認するため，YCrCb輝度補正を使用している．

\section{ちらつき対策}

入力映像の色はフレームごとに変化するため，取得した色をそのまま使用するとモデルの色がちらつく可能性がある．
そこで，前フレームまでの補正値と現在フレームの補正値を混合する平滑化処理を行う．

\begin{lstlisting}[language=Python]
color_match_smoothing = 0.85
\end{lstlisting}

これは，前回までの値を85\%，現在フレームの値を15\%程度反映することを意味する．
これにより，急激な色変化を抑え，安定した表示を実現する．

\section{実装上の処理手順}

処理の流れを以下に示す．

\begin{enumerate}
  \item モデルテクスチャ画像からMediaPipe FaceMeshで顔ランドマークを検出する
  \item 肌色基準ランドマーク群の周辺RGB値を取得し，中央値を $\mathbf{C}_{model}$ とする
  \item 入力映像の各フレームから同じランドマーク群の周辺RGB値を取得し，中央値を $\mathbf{C}_{input}$ とする
  \item 平滑化により，急激な色変化を抑える
  \item 補正式をモデルテクスチャ全体に適用し，OpenGLテクスチャを更新する
\end{enumerate}

\section{確認用のランドマーク描画}

補正に使用しているランドマークを確認できるように，入力映像上へ肌色基準ランドマークを描画している．
描画はオレンジ色の点，白い外枠，ランドマーク番号で行う．
これにより，どの顔領域を肌色基準として使用しているかを視覚的に確認できる．

\section{利点}

本手法の利点は以下である．

\begin{itemize}
  \item 入力映像の肌色に合わせてモデルの色味を自動補正できる
  \item 複数ランドマークの中央値を用いるため，1点基準より安定する
  \item 元テクスチャとの差分を残すため，口や鼻などの情報が消えにくい
  \item RGB補正と輝度補正を切り替えることで，色味合わせと露出合わせを比較できる
\end{itemize}

\section{課題}

一方で，以下の課題が残る．

\begin{itemize}
  \item 顔の左右で照明差が大きい場合，モデル内の元テクスチャの陰影差も残る
  \item 肌色基準ランドマークに髪，眉，影が重なると補正色がずれる可能性がある
  \item RGB差分補正では色味を強く合わせられる一方，色かぶりの影響でモデル全体の色相が揺れる可能性がある
  \item 輝度補正では色味を保てる一方，モデル肌色と入力肌色の差が大きい場合は色味のずれが残る
\end{itemize}

今後の改善案として，顔領域内の部位ごとに補正量を変える方法や，RGB差分補正と輝度補正を混合する方法が考えられる．

\section{まとめ}

複数の顔ランドマークを用いた肌色中央値に基づき，モデルテクスチャの色味補正を実装した．
本手法では，入力映像とモデルテクスチャの肌色基準を比較し，元テクスチャの相対的な色差を保持したまま，モデル全体の色を入力映像に近づける．
また，RGB差分補正，YCrCb輝度補正，Lab輝度補正を切り替えられるようにすることで，肌色そのものを合わせる場合と，モデル本来の色味を保ったまま明るさだけを合わせる場合を比較できる．
これにより，単一ランドマーク基準より安定した色補正が可能となる．

\end{document}
